\section{Introduction}
\label{sec:introduction}

\IEEEPARstart{L}{LM}-based agentic systems have moved from single-shot answer generators to delegated collaborators that perform consequential work --- drafting research papers, reviewing code, interpreting results, and improving proposals~\cite{taskaware2026, mast2025}. Multi-agent variants extend this to teams of LLMs that decompose tasks, parallelize execution, and integrate diverse competencies~\cite{whenagree2024}. Yet collaborator-like behaviour does not by itself produce reliable collaboration: the dominant orchestration paradigms --- routing, aggregation, post-hoc failure analysis --- focus on either selecting an agent before or selecting an output after, and rarely interrogate what an agent has \emph{committed to} between those two points. Recent work shows that interpretation quality~\cite{consistency2024}, ambiguity in user requests~\cite{clamber2024, lhaw2024}, miscalibrated uncertainty~\cite{metacog2024, confidencedichotomy2024}, and behavioural variance under identical prompts~\cite{sameprompt2024} all matter for whether the executed work meets the user's actual need. What unifies these observations is that something happens between the natural-language delegation and the agent's first execution step --- and the existing pipeline rarely exposes it.

A central source of this opacity is that natural-language delegations are usually \emph{operationally underspecified}. A request like ``Improve this paper draft.'' does not specify whether the agent should perform conceptual reconstruction, evidence-claim alignment, novelty positioning, or sentence polish; nor does it specify in what proportion. Several internally consistent executions of the same delegation lead to outputs that the user would judge differently. We call the agent's pre-execution choice across these alternatives a \emph{responsibility projection}: a structured commitment, before any artifact is touched, to a particular subset of operational responsibilities. When the projection differs from the user's later-preferred responsibility structure, we have a \emph{projection mismatch}. Existing routing decisions select an agent before this commitment; existing aggregation steps select among outputs after; existing post-hoc taxonomies label failures that the commitment may have caused. None of these pipelines makes the commitment itself measurable.

This paper makes the pre-execution commitment a first-class object. We formalize the projection $\pi(d, u, t, a) \to r$ as a multi-label weight prediction over a closed responsibility space $J_{v1.1}$ of $|J| = 12$ dimensions, partitioned into seven category dimensions (R1.1--R1.7) for paper-research delegation and five cross-cutting dimensions (RX.1--RX.5) that apply to any delegated work. We use the closure to make projections from different model families directly comparable, and we use that comparability to define and measure projection mismatch as cosine and Jaccard distance between projections. We then propose a closed-loop protocol --- responsibility-bearing delegation --- that projects $r$, solicits responsibility-bearing bids from candidate agents, triggers clarification when projection uncertainty is high relative to clarification cost, selects an agent under a responsibility-adjusted utility, executes, settles per dimension, and updates per-dimension reputation. The protocol is described in Section~\ref{sec:protocol} and stated as Algorithm~\ref{alg:rbd}.

We instantiate the taxonomy and evaluate the projection layer (with a pilot projection-injection treatment for execution) on a 50-example dataset, \textsc{AmbiguousDelegation-50-R1}, constructed from $32$ synthetic and $18$ modified-real artifacts under the closed taxonomy. The empirical contributions are: (i) projection mismatch is a measurable, family-attributable quantity at $T = 0$, exceeding within-family stochastic variance at $T = 0.5$ by an order of magnitude (median $R = 14.4$, paired bootstrap CI excludes zero, Wilcoxon $p < 0.0001$); (ii) under a redesigned 12-judge Anthropic-excluded panel and a three-condition execution split, projection-driven execution shows a directional \emph{disadvantage} relative to direct execution on the headline weighted-R1 settlement loss, with the cost concentrated on R1.4 (novelty mapping) and R1.7 (citation audit) --- the two dimensions whose $s = 5$ anchor demands deep specialty engagement; (iii) a closed Stage 1 human-anchor pilot on R1.7 independently surfaces a methodological boundary condition: form-embedded sharpened anchors are insufficient for reliable rater application without a separately-read protocol document, and the human-anchor scalability constraint is anchor specifiability and domain-expertise gating, not rater throughput. The two findings converge on R1.7 directly (Stage 1 provides a direct boundary-condition signal on R1.7); R1.4 is inferential, since R1.4 was deferred from the Stage 1 closure.

\subsection*{Contributions}

\begin{itemize}
    \item \textbf{Problem formulation.} We formalize multi-agent LLM orchestration as incomplete delegation under natural-language ambiguity, and we identify projection mismatch as a pre-execution failure mode that is structurally distinct from routing failure, aggregation failure, miscalibration, and post-hoc misclassification (Section~\ref{sec:problem}).
    \item \textbf{Closed responsibility taxonomy.} We instantiate a closed 12-dimension responsibility taxonomy $J_{v1.1}$ for paper-research delegation; the closure is what makes pre-execution responsibility projections from different model families directly comparable. We additionally specify a closed-loop protocol (project $\to$ bid $\to$ clarify $\to$ select $\to$ execute $\to$ settle $\to$ reputation) for completeness (Section~\ref{sec:protocol}); this paper evaluates only the projection layer of that protocol.
    \item \textbf{P1 --- Measurability (primary).} On 50 controlled examples, three distinct model families produce projections that differ from each other by an order of magnitude more than five repeat runs of any single family at non-zero temperature ($R(d)$ median $14.4$); all five guard hypotheses pass at $p < 0.0001$. Projection mismatch is measurable and family-attributable (Sections~\ref{sec:results:exp1}, \ref{sec:results:exp1b}).
    \item \textbf{P2 --- Boundary condition on actionability (negative result, reported transparently).} Under a 12-judge Anthropic-excluded panel and a three-condition execution split, projection-driven execution scores worse than direct execution on the headline weighted-R1 settlement loss; the cost concentrates on R1.4 and R1.7. We propose one consistent explanation (broad-attention dilution on dim-specific deep work) and discuss two alternative explanations (format-injection, hard-dim null) that the pilot cannot fully separate (Sections~\ref{sec:results:exp2}, \ref{sec:results:mechanism}).
    \item \textbf{P3 --- Anchor specifiability ceiling (methodological).} A closed Stage 1 human-anchor pilot on R1.7 surfaces a methodological boundary condition that any project proposing crowdworker human anchor as the validity backstop for an LLM-only annotator panel must address: form-embedded sharpened anchors are insufficient for reliable rater application without a separately-read protocol document. The convergence with P2 is direct on R1.7 (where Stage 1 applies) and inferential on R1.4 (not in the Stage 1 closure) (Sections~\ref{sec:results:limitations:specifiability}, \ref{sec:discussion:specifiability}).
\end{itemize}

\textbf{Scope of this paper.} We instantiate the protocol on a single delegation category (R1, paper-research) with five cross-cutting dimensions, and we evaluate at pilot scale ($n = 50$) with one author and two LLM annotators on the hidden-intent track. The other delegation categories in our long-term plan (code review, result interpretation, proposal improvement, technical summarization), the full main-run sample size of $n = 300$, the full four-condition Experiment 2 with task-aware-routing and CLAMBER-style baselines, and the longitudinal reputation experiment on real LLM agents are reserved for follow-up work; their omission is disclosed throughout. The pilot's planned human-annotator extension was attempted, partially executed, and yielded a methodological boundary-condition finding on anchor specifiability, reported in Sections~\ref{sec:results:limitations:specifiability} and \ref{sec:discussion:specifiability}. The contribution of the present paper is that responsibility projection becomes \emph{measurable} at this scope, naive projection injection is not yet actionable in current form, and human-anchor scalability is bound by anchor specifiability and domain-expertise gating; main-run extensions are the natural next step.

The rest of the paper proceeds as follows. Section~\ref{sec:related} positions the work against adjacent orchestration, ambiguity, calibration, evaluation, and annotation frames. Section~\ref{sec:problem} fixes the closed taxonomy and the projection-mismatch quantity. Section~\ref{sec:protocol} states the protocol and Algorithm~\ref{alg:rbd}. Sections~\ref{sec:experiments}--\ref{sec:results} report the pilot setup and findings. Section~\ref{sec:discussion} discusses implications and Section~\ref{sec:conclusion} concludes.
